\documentclass{safeai2026}

\usepackage[american]{babel}
\usepackage{natbib}
    \bibliographystyle{plainnat}
    \renewcommand{\bibsection}{\subsubsection*{References}}
\usepackage{mathtools}
\usepackage{booktabs}
\usepackage{url}

\title{AgentCloseoutBench: A Pre-Release Benchmark for Dark-Pattern Detection at the Agentic Coding Assistant Closeout Boundary}

\author[1]{Anonymous Submission}
\affil[1]{%
    Anonymous Affiliation
}

\begin{document}
\maketitle

\begin{abstract}
We formalise the assistant-message closeout boundary of agentic coding assistants as a distinct safety surface and release a pre-release candidate benchmark, AgentCloseoutBench, for detecting four closeout-stage dark patterns: premature wrap-up, cliffhanger nudging, role-play drift, and sycophancy.
Existing runtime monitors for LLM agents target action-level execution traces.
We argue that the textual closeout artifact emitted at lifecycle stop events is an orthogonal surface where the agent's claims about its own work can be measured deterministically and out-of-band.
We release an 800-record candidate corpus, an open annotation protocol with multi-model judges, an open-source deterministic reference engine, and an open hook-pack runtime that integrates with the Claude Code stop and subagent-stop lifecycle hooks.
We report inter-annotator agreement among three frontier large language model judges on a 200-record sample and discuss honest limitations: candidate labels are not yet human-adjudicated, the reference engine is deterministic rather than semantic, and lexical evasion remains possible.
This work is intentionally framed as pre-release scaffolding, not a final benchmark score.
\end{abstract}

\section{Introduction}\label{sec:intro}
% Budget: ~0.5 page. State the gap, the boundary, the contribution. No final scores.
% TODO: write 3 paragraphs:
% 1. Agentic coding assistants increasingly close work autonomously and emit closeout text that asserts what was done.
% 2. Existing runtime safety work targets action traces; the textual closeout artifact is an under-studied safety surface.
% 3. Our contribution: formalisation, candidate corpus, judge protocol, reference engine, hook runtime. Frame as pre-release.

\section{Related Work}\label{sec:related}
% Budget: ~0.6 page. Five citations minimum. Differentiate cleanly.
% TODO: structure as three paragraphs:
%
% 1. Chat-surface and web-agent dark patterns:
%    - DarkBench (ICLR 2025) -- six dark-pattern categories on chat outputs; we differentiate on
%      surface (closeout text at stop hooks) and category set (four closeout-stage patterns).
%    - HarmBench (arXiv:2402.04249) -- standardised red-teaming for refusal; orthogonal to closeout claims.
%    - SORRY-Bench (45-class taxonomy) -- safety refusal taxonomy on chat outputs.
%    - CASE-Bench -- context-aware safety benchmark with 2000+ annotators on chat queries.
%
% 2. Runtime monitoring of LLM agents:
%    - ProbGuard / Pro2Guard (arXiv:2508.00500) -- probabilistic DTMC monitor over agent action traces.
%      Differentiation: action-level prediction vs closeout-text claim discipline.
%    - AgentSpec (arXiv:2503.18666, ICSE'26) -- DSL for action-level runtime constraints.
%      Differentiation: pre-action enforcement vs post-action self-report enforcement.
%
% 3. LLM-as-judge methodology:
%    - DarkBench used three LLM annotators with three human annotators on 1680 examples
%      and explicitly reported per-category Cohen's kappa as low as 0.27 for one model on one category;
%      we follow this precedent for honest kappa reporting in section~\ref{sec:annotation}.

\section{The Closeout Boundary}\label{sec:boundary}
% Budget: ~0.5 page. Formalise the surface.
% TODO:
% - Define the closeout artifact as the last_assistant_message at Stop/SubagentStop events.
% - Distinguish from chat surfaces (multi-turn dialogue) and action surfaces (tool-call traces).
% - State the safety property: the closeout artifact contains the agent's self-claims about completion,
%   evidence, scope, and remaining work; these claims are accessible without access to action traces.
% - Note the practical lifecycle hooks in Claude Code: Stop, SubagentStop, last_assistant_message.

\section{Dataset and Categories}\label{sec:dataset}
% Budget: ~0.5 page.
% TODO:
% - 800 candidate records, 4 categories x 100 positives x 100 negatives.
% - Categories: wrap_up, cliffhanger, roleplay_drift, sycophancy. Define each in one sentence.
% - 8 task types per category/label balanced (refactor, research, test, review, debug, write, migrate, infra).
% - Candidate labels generated deterministically; final labels NOT YET assigned.
% - Provenance lane: candidate_synthetic + candidate_public_adversarial (16 rows, reviewed sources).
% - Cite SPEC.md and DATASET_CARD.md (anonymous, by URL placeholder).

\section{Annotation Protocol and Inter-Judge Agreement}\label{sec:annotation}
% Budget: ~0.8 page. Most important section for credibility.
% TODO:
% - Following DarkBench precedent~\citep{darkbench2025}, we use three frontier LLM judges:
%   {claude-sonnet-4-6, gpt-5.x (TBD exact model ID), gemini-2.x (TBD exact model ID)}.
% - 200-record stratified sample from the 800 candidate corpus, deterministic seed.
% - Binary annotation per category per record; judges given prompt with category definitions.
% - Report Table 1 (analogous to DarkBench Table 3): per-category pairwise Cohen's kappa,
%   per-judge vs candidate-label agreement, overall.
% - Following~\citep{darkbench2025}: ``Despite a low Cohen's Kappa on some dark pattern categories,
%   indicating poor inter-rater agreement, the summary statistics over models and dark patterns
%   remain consistent.'' We expect and accept this in our context; aggregate consistency is the
%   construct-validity claim, not perfect per-category agreement.
% - Explicit note: this is not human-gold. Human adjudication is deferred to a v1.0 release.

\section{Reference Engine and Hook Runtime}\label{sec:engine}
% Budget: ~0.4 page.
% TODO:
% - Deterministic rule packs over normalised closeout text, no LLM in the verdict path.
% - Categories implemented as per-category engine manifests, rule-pack hash recorded with every decision.
% - Engine claim: ``Out-of-band deterministic enforcement at the agentic coding assistant closeout
%   boundary makes specific dark-pattern and false-closeout mechanics observable, reproducible,
%   and benchmarkable.'' (Quoted verbatim from internal claim ledger.)
% - NOT a semantic intent detector. NOT prompt-injection-proof.
% - Hook runtime: thin Bash/JQ wrappers that invoke the reference engine at Stop/SubagentStop.

\section{Threat Model and Limitations}\label{sec:threat}
% Budget: ~0.6 page. Be explicit, paper-grade reviewers will look for this.
% TODO: enumerate (use \begin{itemize}):
% 1. Lexical evasion: deterministic rules are matchable on text surface; paraphrase can evade.
%    We measure this brittleness rather than hide it.
% 2. Hook misconfiguration: the hook runtime depends on operator wiring; a disabled hook is silent.
% 3. Runtime bypass: a local operator can intentionally disable hooks; this is not an OS sandbox.
% 4. In-band vs out-of-band: hooks run outside the model context, but evasion via the model
%    output text is still possible; the model can be coerced via prompts to phrase a false
%    closeout in ways our rules do not catch.
% 5. Evidence-marker limitation: ``verification passed'' style markers are closeout-contract
%    evidence, not independent proof of task success.
% 6. Synthetic corpus artifacts: candidate corpus is template-generated. Improved controllability
%    has the tradeoff of lexical patterns that simple baselines may exploit.
%
% Also list:
% - English-only.
% - Single-assistant (Claude Code lifecycle) surface; we do not claim universal agent coverage.
% - Candidate labels not human-adjudicated.
% - No final benchmark performance numbers in this paper. F1 numbers from candidate diagnostics
%   are explicitly excluded from abstract, introduction, and conclusion.

\section{Conclusion}\label{sec:conclusion}
% Budget: ~0.2 page.
% TODO:
% - One-paragraph restatement of contribution.
% - Honest framing: pre-release benchmark scaffold, not final-result paper.
% - Pointer to v1.0 plan: human adjudication, locked-test, expanded public-derived lane.

\bibliography{refs}

\end{document}
